% Source: https://artofproblemsolving.com/community/c3h2006585_kootrapalis_introductory_proof_problems_kipp
% Stable identity: aops-topic-2006585-q05
Let $\Gamma$ be the circumcircle and $\ell$ the line $AO$. Let $\sigma$ be the reflection in $\ell$.

\begin{asy}
usepackage("amsmath");
size(4.5cm);
defaultpen(fontsize(10pt));
pair O=(0,0); real R=2;
pair A=R*dir(100), B=R*dir(215), C=R*dir(-15);
draw(circle(O,R),gray(0.5));
draw(A--B--C--cycle,linewidth(0.9));
pair Ap=-A;
draw((1.25*A)--(1.15*Ap),blue+dashed);
dot(O); dot(A); dot(B); dot(C);
label("$A$",A,N); label("$B$",B,SW); label("$C$",C,SE); label("$O$",O,E);
label("$\ell$",1.15*Ap,S,blue);
label("\mbox{}",(min(currentpicture,user=true)+max(currentpicture,user=true))/2,black); // keeps the text after the figure black when LaTeX typesets the labels inline
\end{asy}

\begin{itemize}
\item Since $\ell$ passes through the centre $O$, $\sigma$ maps $\Gamma$ onto itself.
\item Since $\ell$ bisects the angle between the lines $AB$ and $AC$ and passes through their common point $A$, $\sigma$ maps line $AB$ onto line $AC$. (Lines $AB$ and $AC$ are different from $\ell$, since otherwise $A,B,C$ would be collinear.)
\end{itemize}
Line $AB$ meets $\Gamma$ in exactly two points, $A$ and $B$. Its image, line $AC$, meets $\Gamma$ in $A$ and $C$. As $\sigma$ fixes $A$ and maps $\Gamma\cap AB$ onto $\Gamma\cap AC$, it must map $B$ to $C$. Reflections preserve distances and fix $A$, so
\[
AB=\sigma(A)\sigma(B)=AC .
\]
